% Modernised reading — fonds Alexandre Grothendieck, shelfmark 127
% (pages 1-23, the whole folder).
%
% This is an INTERPRETATION, not a transcription. It re-reads the pages in
% today's notation and vocabulary, states the hypotheses the manuscript leaves
% implicit, and drops the critical apparatus so that the mathematics reads
% straight through. Everything the brackets and underlines carried moves into a
% footnote. It is held to being mathematically correct as it stands; where the
% manuscript is loose or mistaken, the true statement is what appears here and
% the footnote says what the page has. In any dispute of substance the
% transcription governs: whoever cites Grothendieck must cite batch-01.fr.tex
% (pp. 1-20) or batch-02.fr.tex (pp. 21-23).
%
% Pass: Opus 5.5 (claude-opus-5-5), 2026-10-03 — first pass, unchecked against the pages by a human.
% The folder was read whole, its two batches together; this is one document
% for the shelfmark. It was made on Opus 5.5 and has not been measured against
% the reference reading of folder 115 (made on Opus 5). The EGA references
% are the page's own where the footnotes say so; those we add are cited from
% memory and were not checked against the sources.
%
% Scope. Page 1 and page 21 are covers in his hand whose words head sections
% here. Page 8 is a typed leaf of the EGA IV typescript (n° 21.10) carrying
% his ink corrections: published text, identified and not restated. Page 16 is
% in another hand: a fair copy of Th. 1.4 (« Lemme 1.4' »); it is used only
% as far as the transcription's note describes it (RIGHTS.md). Much of the
% connective prose of pp. 3-4, 7, 9 and 13 is illegible in the transcription;
% the reading gives what carries and says where it reconstructs.
%
% Spine. Two runs under two covers, plus scratch.
%   pp. 1-6, 14-16, 18-19  cover « Application de la descente non plate » ;
%       § 1 « Platitude d'une adhérence schématique »: problem 1.1 (flat
%       quotient G of F extending G_U, U universally G-dense rel. Y), 1.2
%       (fpqc-local), 1.3 (non-flat descent over an artinian base along an
%       epimorphism), Th. 1.4 (conditions (a) density, (b) solutions over a
%       separating family Y_i -> Spec Ô_{Y,y}, (c) finiteness), proof outline
%       (p. 4), Cor. 1.5 (struck), 1.6-1.7 (valuative criterion). Pencil
%       drafts pp. 14-15, 18-19 (Remarques 1.4.2-1.4.3, « Vrai théorème »,
%       Chevalley-type finite subfamily, p. 19).
%   pp. 7, 9, 10, 12  Th. 1.8: Hironaka's lemma -> openness of the locus
%       where X -> Y is flat with geometrically normal fibres; proof sketch
%       reducing flatness to the valuative criterion of § 1.
%   p. 8   typed EGA IV leaf.
%   pp. 11, 13, 17, 20  the site of finite schemes; quasi-coherence and its
%       failure for kernels; « faisceaux formels »; the Milnor lim^1 sequence.
%   pp. 21-23  cover « Foncteurs sur morphismes non ramifiés, étales … »:
%       Lemma (sheaf for effective epis + reduction-injective => exact for
%       all epis), Corollary (surjective => injective), Proposition on
%       abelian subschemes; struck Corollary breaks off.
%
% Notation fixed for the document: Y the base throughout (the page writes S
% at pp. 2 and 15, and from p. 11 on S is a different base; kept apart);
% i : U -> X; G := Im(F -> i_* G_U); « solution » = solution of Pb 1.1;
% T = closure of a point x' of Ass; C the site of finite schemes (p. 11, 13,
% 17) and C the category of finite S-schemes (p. 22) are both his round C
% and are said apart where they meet.
%
% Substantive departures, each footnoted where it occurs:
%  p. 2   « univ. injectif rel/Y » read as: injective after every base change
%         Y' -> Y (EGA IV 11.10 vocabulary, ours); U retrocompact needed for
%         i_* quasi-coherent (the page's pencil addition, uncertain reading).
%  p. 3   (a) restated in the density form T_y cap U dense in T_y; the
%         equivalence with the page's specialisation form is ours.
%  p. 4   proof given as an outline; the link (b) + 1.3 => solutions over
%         every Spec(Ô/m^n), via a finite subfamily (p. 19), is ours.
%  p. 5   Cor. 1.5 struck on the page; reported, not asserted.
%  p. 14  M -> N with M_alpha ≅ N_alpha: the page concludes only at primes
%         not containing J; over a J-adically complete A, M -> N is an iso.
%  p. 17  exponent of lim^1 is q-1 (the page reads q).
%  p. 20  lim^(p) = 0 for p >= 2 needs a countable index (supplied).
%  p. 22  the missing injectivity of F(X' x_Y X') -> F(X x_Y X) supplied
%         from p. 23's corollary (ours); the factorisation into effective
%         epimorphisms is assumed by the page, not proved.
%  p. 23  S' finite over S supplied; « abélien » a doubtful reading.
%
% Announced and never established: Lemma 1.3 (no proof); the proof of 1.4
% (p. 4 largely illegible); 1.5-1.7 (no proofs); the « Vrai théorème »
% (p. 15); Th. 1.8 beyond the sketch; the « Lemme à dégager » (p. 9); the
% comparison H^*(C_/delta, F) ≅ H^*(Ẑ, F) (p. 13, marked « à vérifier »);
% « Dans le cas qui nous occupe » (p. 20, stops); the factorisation of
% epimorphisms (p. 22); the struck Corollary of p. 23.
\documentclass[11pt,a4paper]{article}
\input{../preamble/grothendieck}

\folder{127}
\pages{1}{23}
\dating{s.d. — le groupe « Descentes » (126 à 133) est daté [avant 1970]}
\watermark{Édition de démonstration\\interprétation personnelle de l'œuvre}
\foldertitle{Descente non plate : notes manuscrites (s.d.) — lecture modernisée du dossier entier}

\begin{document}

\section*{Résumé}

\begin{resume}
Ces vingt-trois feuillets cherchent à quelle condition une famille d'objets
géométriques qui « varie bien » sur une partie de l'espace continue à varier
bien quand on la prolonge à tout l'espace.

Le cadre est celui d'une famille d'espaces $X$ paramétrée par une base $Y$,
et d'un faisceau de modules $F$ sur $X$. Sur un ouvert $U$ de $X$, on s'est
donné un quotient de $F$ qui est \emph{plat} sur $Y$ : en termes imagés, ses
fibres au-dessus des points de la base forment une famille sans sauts, où rien
n'apparaît ni ne disparaît brusquement quand le paramètre bouge. Le
problème est de prolonger ce quotient à tout $X$ en gardant la platitude, et
de façon que ce qui se passe hors de $U$ soit entièrement commandé par ce qui
se passe sur $U$. Une analogie aide : une suite de courbes du plan qui
converge ; la courbe limite est imposée, c'est l'« adhérence », et la seule
question est de savoir si elle a le même degré que les autres, ou si elle a
ramassé en route un morceau parasite. Il n'y a qu'un candidat ; tout est de
savoir s'il convient.

L'idée de Grothendieck est de juger le candidat sur des bases très petites.
On ne peut pas tester la platitude point par point, mais on peut la tester
après des changements de base vers des espaces minuscules : des épaisseurs
infinitésimales d'un point (les schémas artiniens) et des germes de courbes
lisses (les « traits »). Le titre du dossier dit l'outil : la \emph{descente
non plate}. La descente ordinaire permet de redescendre une construction faite
sur un revêtement fidèlement plat ; ici les morphismes le long desquels on
redescend ne sont pas plats, seulement des épimorphismes, et le prix à payer
est de ne redescendre que des quotients plats d'un faisceau déjà donné. Le
théorème central (pages 3 et 4), sous la forme qu'un brouillon de la page 15
appelle le « vrai théorème », dit en substance : sous des hypothèses de
finitude et de régularité des complétions, le prolongement existe si et
seulement s'il existe après tout changement de base vers une épaisseur de
point ou vers un trait. Le dossier n'en écrit pas de preuve complète. C'est le modèle de ce qu'on appelle le critère
valuatif de platitude.

La deuxième partie applique ce critère à un résultat de géométrie classique,
le lemme de Hironaka : si une famille de variétés normales dégénère en une
fibre dont la partie réduite est encore normale, sans multiplicité aux points
génériques, alors la fibre est elle-même normale et la famille est plate près
d'elle. Le dossier en tire un énoncé d'ouverture (page 7, remise au net
page 10), et en esquisse la preuve : on se ramène à un trait, et la platitude
vient de la première partie.

Viennent ensuite des pages plus détachées : un site formé des schémas finis
sur un corps, où l'on voit qu'un noyau n'est pas quasi-cohérent faute de
platitude ; des « faisceaux formels » et la suite exacte qui mesure ce qu'on
perd en passant à la limite projective ; enfin, sous une seconde chemise, un
lemme de descente pour des foncteurs sur les schémas finis, qui fait passer
la condition de faisceau des épimorphismes effectifs à tous les
épimorphismes, et son application aux sous-schémas abéliens. C'est, en
germe, l'outil dont la première partie a besoin.

Un feuillet dactylographié d'EGA IV, corrigé à la main, s'est glissé au
milieu ; une page d'une autre main met au net le théorème central. Rien ne
date ces feuilles ; l'inventaire range le groupe avant 1970. Pour aller plus
loin : adhérence schématique, densité schématique universelle (EGA IV, §
11.10), critère valuatif de platitude, lemme de Hironaka, fibres formelles,
épimorphismes effectifs.

\keywords{schematic closure, flatness, non-flat descent, universally schematically dense open, valuative criterion of flatness, artinian base change, formal fibres, henselization, Chevalley's lemma, Hironaka's lemma, normal morphism, universally open morphism, site of finite schemes, quasi-coherent sheaf, formal sheaf, Milnor exact sequence, derived inverse limit, effective epimorphism, rigidity of abelian schemes}
\end{resume}

\pagerange{1}{23}
\section*{Le fil du dossier, et les conventions}

\subsection*{Les stations}

Le dossier a deux chemises de sa main, pages 1 et 21, et se lit en cinq
suites.

\begin{enumerate}
  \item \emph{Pages 1 à 6, avec les brouillons des pages 14, 15, 18 et 19.}
  La chemise « Application de la descente non plate : platitude d'une
  adhérence schématique, et lemme de Hironaka », puis le § 1 : le problème
  1.1, sa nature locale (1.2), la descente non plate sur une base artinienne
  (1.3), le théorème 1.4 et l'esquisse de sa preuve, les corollaires 1.5 à
  1.7. La page 16, d'une autre main, met au net l'énoncé 1.4.
  \item \emph{Pages 7, 9, 10 et 12.} Le théorème 1.8, application du lemme de
  Hironaka, sa démonstration esquissée, une remise au net de l'énoncé
  (page 10) et la suite de la preuve (page 12).
  \item \emph{Page 8.} Un feuillet du tapuscrit d'EGA IV, sans rapport avec
  le reste.
  \item \emph{Pages 11, 13, 17 et 20.} Le site des schémas finis,
  quasi-cohérence, faisceaux formels, suite exacte de $\varprojlim^{1}$.
  \item \emph{Pages 21 à 23.} La chemise « Foncteurs sur morphismes non
  ramifiés, étales … », un lemme de descente pour les foncteurs sur les
  schémas finis, son corollaire, et le cas des sous-schémas abéliens.
\end{enumerate}

Le lien entre la première et la dernière suite est de nous, mais il est
étroit : le lemme 1.3 redescend un objet le long d'un épimorphisme non plat
de schémas artiniens, et le lemme de la page 22 dit exactement quand un
foncteur sur les schémas finis vérifie la condition de faisceau pour
\emph{tous} les épimorphismes. Le dossier ne relie pas lui-même les deux.

\subsection*{Les conventions}

\begin{itemize}
  \item La base s'appelle $Y$ dans toute la première partie ; les pages 2 et
  15 écrivent aussi $S$ pour elle\footnote{À la page 2, « univ.\ injectif
  rel/$S$ » et « de prés.\ finie sur $S$ » ; à la page 15, les changements de
  base $S' \to S$. Rien ne distingue là $S$ de $Y$.}. À partir de la page
  11, $S$ est une autre base (le spectre d'un corps, puis un préschéma
  localement noethérien).
  \item $f \colon X \to Y$, $U \subset X$ ouvert, $i \colon U \to X$
  l'inclusion ; pour un point $y \in Y$, $X_y$, $U_y$ sont les fibres. Pour
  un module $G$, $G_y$ est sa restriction à la fibre $X_y$ et
  $\mathrm{Ass}\, G_y$ l'ensemble de ses points associés.
  \item On note toujours $G := \mathrm{Im}\bigl(F \to i_*(G_U)\bigr)$ le
  candidat du problème 1.1, et l'on appelle \emph{solution} une solution de
  ce problème.
  \item Un \emph{trait} est le spectre d'un anneau de valuation discrète.
  \item Les deux catégories que la page note $\mathcal{C}$ sont distinctes :
  à la page 11, un site de schémas finis sur un corps ; à la page 22, la
  catégorie des préschémas finis sur une base $S$ localement noethérienne.
\end{itemize}

\subsection*{Ce que le dossier annonce et n'établit pas}

Le lemme 1.3 n'a pas de démonstration ; celle du théorème 1.4 n'est lisible
qu'en partie ; les corollaires 1.5 à 1.7 et le « Vrai théorème » de la page
15 sont énoncés sans preuve ; le théorème 1.8 n'est démontré qu'en esquisse,
avec un « Lemme à dégager » jamais dégagé ; la comparaison cohomologique de
la page 13 est marquée « à vérifier » ; la page 20 s'arrête sur « Dans le cas
qui nous occupe, il vient : » ; le corollaire biffé de la page 23 s'arrête
sur son énoncé.

\pagerange{1}{2}
\section*{I. Platitude d'une adhérence schématique}

La chemise porte, soulignés, les mots « Application de la descente non
plate : platitude d'une adhérence schématique, et lemme de
Hironaka »\footnote{Page 1, feuillet de garde ; le mot lu « lemme » est
repassé ou surchargé.}. Ce que la page 2 ouvre est le § 1, que la page
intitule « Platitude d'une adhérence schématique ».

\subsection*{Le problème 1.1 (page 2)}

Soient $f \colon X \to Y$ un morphisme de schémas, $U$ un ouvert
rétrocompact de $X$\footnote{« rétrocompact » est ajouté au crayon, encerclé,
et sa lecture est peu sûre. L'hypothèse est nécessaire : elle rend
$i \colon U \to X$ quasi-compact, donc $i_*$ conserve la quasi-cohérence, et
c'est ce qui fait de $G$ ci-dessous un module quasi-cohérent.}, $F$ un
$\mathcal{O}_X$-module quasi-cohérent et $G_U$ un quotient quasi-cohérent de
$F|_U$, \emph{plat} sur $Y$.

\emph{Problème.} Trouver un quotient quasi-cohérent $G$ de $F$ qui prolonge
$G_U$, qui soit plat sur $Y$, et tel que le morphisme $G \to i_*(G_U)$ soit
\emph{universellement injectif relativement à $Y$} : pour tout changement de
base $Y' \to Y$, de notations $X' = X \times_Y Y'$, $U' = U \times_Y Y'$,
$i' \colon U' \to X'$, le morphisme $G_{Y'} \to i'_*(G_{U'})$ est
injectif\footnote{La page écrit « univ.\ injectif rel/$S$ » et, au-dessus,
« [$G \to i_*(G_U)$ » ; la formulation par changement de base est la nôtre.
Dans le vocabulaire d'EGA IV, § 11.10, elle dit que $U$ est
\emph{universellement $G$-schématiquement dense relativement à $Y$}.}.

\emph{Unicité.} Si le problème a une solution, c'est
\[
  G = \mathrm{Im}\bigl(F \to i_*(G_U)\bigr) ,
\]
puisque $F \to G$ est surjectif et $G \to i_*(G_U)$ injectif. La question est
donc seulement de savoir si ce candidat est plat sur $Y$ et universellement
injectif dans $i_*(G_U)$ : c'est la platitude d'une \emph{adhérence
schématique}, d'où le titre.

Une remarque entre crochets donne le critère commode\footnote{La page écrit
« Il faux et suffit » (deux mots au crayon, peu sûrs) : nous lisons « il faut
et il suffit ». L'énoncé exact que nous donnons est celui d'EGA IV,
11.10.9--11.10.10, cité de mémoire ; la page n'y renvoie pas.} : si $f$ est
localement de présentation finie et $G$ de présentation finie et plat sur
$Y$, la condition d'injectivité universelle équivaut à
\[
  \mathrm{Ass}\, G_y \subset U_y \qquad \text{pour tout } y \in Y .
\]
Une note de marge, barrée et en partie illisible, envisage de se ramener à
$X$ plat sur $Y$ et à $F$ de type fini\footnote{Note écrite en long dans la
marge gauche, barrée de traits obliques, au début illisible.}.

\subsection*{Nature locale (lemme 1.2)}

Le problème est local sur $Y$ pour la topologie fpqc : il a une solution
s'il en a une après un changement de base fidèlement plat et
quasi-compact. Il en va de même de la variante où l'on demande en outre $G$
localement de présentation finie. Plus généralement, on peut redescendre le
long d'un morphisme fidèlement plat quasi-compact $X' \to X$\footnote{La marge
renvoie à SGA 1, VIII (descente fidèlement plate). Justification, de nous :
la formation de $i_*$ et celle d'une image commutent au changement de base
plat (c'est encore la rétrocompacité de $U$ qui sert) ; la platitude sur $Y$
et l'injectivité descendent par fidèle platitude, et les changements de base
$Y'' \to Y$ préservent la fidèle platitude de $X' \to X$.}.

\pagerange{2}{3}
\subsection*{La descente non plate sur une base artinienne (lemme 1.3, pages 2 et 3)}

La phrase qui introduit le lemme dit ce qu'est la descente non plate : la
descente « des quotients \emph{plats} d'un faisceau déjà donné ».

\emph{Lemme 1.3.} Supposons $Y = \mathrm{Spec}\, A$ artinien, et soit
$Y' \to Y$ un épimorphisme de schémas tel que le problème ait une solution
$G'$ après le changement de base $Y' \to Y$. Supposons de plus, ou bien
$Y' \to Y$ fini, ou bien $f$ localement de présentation finie et $G'$ de
présentation finie. Alors le problème a une solution sur $Y$ ; dans le second
cas, elle est de présentation finie\footnote{L'énoncé est coupé au bas de la
page 2 (« Suppo\ldots ») et repris en haut de la page 3. L'alternative des
hypothèses vient d'une note au crayon de la page 2 (« Il faut supposer ici
$Y' \to Y$ fini \ldots\ ou alors de prés.\ finie ») et du début de la page
3. Le lemme n'est pas démontré dans le dossier.}.

Le mot « épimorphisme » a ici toute sa force, et c'est ce qui rend
l'énoncé non trivial : quand $Y$ est local artinien, un morphisme
$Y' \to Y$ est un épimorphisme dès qu'il est surjectif et que
$A \to \Gamma(Y', \mathcal{O}_{Y'})$ est injectif\footnote{Vérification de
nous : deux morphismes $Y \rightrightarrows Z$ qui coïncident sur $Y'$
coïncident sur le point de $Y$, donc tombent dans un même ouvert affine de
$Z$, et sont alors déterminés par des homomorphismes d'anneaux à valeurs dans
$A$, que l'injectivité sépare.}. Ainsi la réunion disjointe de deux copies
de $\mathrm{Spec}\, k[\varepsilon]$, envoyées sur les deux axes de
$\mathrm{Spec}\, k[x,y]/(x,y)^2$, est un épimorphisme qui n'est pas plat.

\pagerange{3}{4}
\section*{II. Le théorème 1.4 (pages 3 et 4)}

\subsection*{L'énoncé}

\emph{Théorème 1.4.} Supposons $X$ et $Y$ localement noethériens, $F$ et
$G$ cohérents, et les anneaux locaux de $X$ aux points de $X - U$ à fibres
formelles géométriquement normales\footnote{Le mot qui suit « géom. » est
illisible à la page 3 ; la page 4 (« à fibres formelles géom.\ normales ») et
le corollaire 1.5 le donnent. « en tout pt de $X - U$ » est une note encadrée
de la marge. « Lemme » est biffé sous « Théorème », et « Néc.\ Suff. » est
ajouté au crayon.}. On fait les hypothèses suivantes.

\begin{itemize}
  \item[(a)] \emph{Densité.} Pour toute spécialisation $y$ d'un point $y'$ de
  $Y$, et tout point $x' \in \mathrm{Ass}\,(G_U)_{y'}$, d'adhérence
  $T = \overline{\{x'\}}$, l'ouvert $T_y \cap U$ est dense dans
  $T_y$\footnote{La page énonce (a) en termes de points : pour tout
  $x \in X_y - U_y$ qui spécialise $x'$, il existe un point $x_1 \in U_y$ qui
  spécialise $x'$ et généralise $x$. Les deux formes sont équivalentes (de
  nous) : dans un espace noethérien, un point est dans l'adhérence d'un ouvert
  de $T_y$ si et seulement s'il spécialise un point de cet ouvert. La forme par
  densité est celle des pages 11, 14 et 16. Les ajouts en interligne qui
  précisent les quantificateurs sont serrés et en partie superposés.}.
  \item[(b)] \emph{Solutions formelles.} Pour tout $y \in Y$, il existe une
  famille de morphismes locaux $Y_i \to \mathrm{Spec}\,
  \widehat{\mathcal{O}}_{Y,y}$, les $Y_i$ noethériens locaux, après chacun
  desquels le problème a une solution, et telle que
  \[
    \bigcap_i \mathrm{Ker}\bigl(\widehat{\mathcal{O}}_{Y,y} \to
    \mathcal{O}_{Y_i, y_i}\bigr) = 0 .
  \]
  On peut prendre pour $Y_i$ les $\mathrm{Spec}\bigl(
  \widehat{\mathcal{O}}_{Y,y} / \widehat{\mathfrak{m}}_y^{\,i+1}\bigr)$,
  $i \in \mathbb{N}$, par le théorème d'intersection de Krull.
  \item[(c)] \emph{Finitude.} Ou bien $f$ est localement de présentation
  finie, ou bien les $Y_i \to \mathrm{Spec}\, \widehat{\mathcal{O}}_{Y,y}$
  sont finis\footnote{La fin de (b) et le début de (c) sont très surchargés à
  la page 3, qui semble écrire « $f$ loc.\ de t.\ finie » et « les $Y_i \to
  Y'$ finis ». L'alternative que nous donnons est celle de la mise au net de la
  page 16 et du lemme 1.3, auquel (c) sert à se ramener.}.
\end{itemize}

\emph{Alors} le problème a une solution, qui est un module cohérent plat sur
$Y$.

\emph{Nécessité.} Si $f$ est localement de type fini, (a) et (b) sont
nécessaires à l'existence d'une solution. Pour (b), c'est immédiat : une
solution reste une solution après tout changement de base, c'est le sens du
mot « universellement » dans 1.1. Pour (a), la page renvoie à EGA IV,
12.1.1.5\footnote{Le renvoi « IV 12.1.1.5 » est net ; la phrase qui le suit
est en grande partie illisible. Nous n'avons pas vérifié ce que dit ce numéro.
Le raisonnement esquissé à la page 14 est le suivant : si $G$ est plat sur
$Y$, un point maximal de $T_y$ est associé à $G_y$ (renvoi « EGA IV
12.2 »), donc dans $U_y$ si $G$ est une solution.}.

\subsection*{La démonstration (pages 3 et 4)}

La page 4 est en grande partie illisible ; voici l'ossature qui
porte\footnote{La reconstruction des liaisons est de nous ; les étapes et les
formules sont celles des pages 3 et 4, et la page 19 reprend la réduction par
la descente fpqc, l'hensélisé et le complété.}.

\begin{enumerate}
  \item \emph{Localisation.} On se localise en un point maximal $x$ de
  l'ensemble des points de $X_y - U_y$ ($y = f(x)$) qui sont dans
  $\mathrm{Ass}\, G_y$. On peut alors supposer $X$ local de point fermé $x$, et
  $U = X - \{x\}$.
  \item \emph{Hensélisation.} Par descente fidèlement plate (lemme 1.2), on
  peut supposer $X$ local hensélien ; les hypothèses se conservent, notamment
  celle sur les fibres formelles, qui est stable par hensélisation.
  \item \emph{Complétion.} Soit $B = \mathcal{O}_{X,x}$. Le morphisme
  $\mathrm{Spec}\, \widehat{B} \to \mathrm{Spec}\, B$ est fidèlement plat, et
  ses fibres sont géométriquement intègres\footnote{C'est ici que sert
  l'hypothèse sur les fibres formelles : pour un anneau local hensélien à fibres
  formelles géométriquement normales, les fibres formelles sont
  géométriquement intègres (EGA IV, § 18.9, cité de mémoire). La page dit
  « fibres (géom.) intègres », « (géom.) » ajouté.}. Les cycles associés
  $Z'_\alpha$ de $G' = G \otimes_B \widehat{B}$ sont donc exactement les
  images réciproques des cycles $Z_\alpha$ de $G$, et la condition (a), qui
  dit $Z_\alpha \cap U \neq \emptyset$ là où il faut, passe de $(X, G, U)$ à
  $(X', G', U')$. La condition (b) se conserve aussi.
  \item \emph{Passage à la limite.} Sur la base complète, les solutions sur
  les $Y_n$ donnent, à la limite, un module quotient $H$ de $G$, plat sur $Y$,
  et tel que $G \to H$ soit injectif sur $U$ ; comme $G \to i_*(G_U)$ est
  injectif, $G \to H$ est un isomorphisme, donc $G$ est plat, et $x \notin
  \mathrm{Ass}\, G_y$.
\end{enumerate}

Le rôle exact de (b) et (c) dans l'étape 4 n'est pas lisible à la page 4. Les
pages 14 et 19 permettent de le reconstituer, et nous le faisons ci-dessous.

\pagerange{14}{15}
\section*{III. Les brouillons du théorème 1.4 (pages 14, 15, 16, 18 et 19)}

\subsection*{Deux lemmes d'algèbre locale (page 14)}

La page 14 est une page de notes au crayon. Deux énoncés s'y lisent.

D'abord, pour $A$ noethérien complet pour la topologie $\mathfrak{J}$-adique,
$A = \varprojlim A/\mathfrak{J}^{\alpha+1}$, et $M \to N$ un morphisme de
$A$-modules de type fini qui induit des isomorphismes
$M_\alpha = M/\mathfrak{J}^{\alpha+1} M \simeq N_\alpha$ pour tout $\alpha$ :
alors $M_{\mathfrak{p}} \simeq N_{\mathfrak{p}}$ pour tout premier
$\mathfrak{p}$ ne contenant pas $\mathfrak{J}$. C'est vrai, et même faible :
$M$ et $N$ étant séparés et complets, $M \to N$ est un
isomorphisme\footnote{La conclusion de la page est la version localisée ; la
remarque qu'on a l'isomorphisme global est de nous. Une parenthèse peu lisible
sous la conclusion mentionne $\widehat{M}_{\mathfrak{p}}$.}.

Ensuite, la nécessité de (a) et un contre-exemple. La page raisonne ainsi : si $G$ est une solution, $x_\alpha \in
\mathrm{Ass}\, G_U$ et $T_\alpha = \overline{\{x_\alpha\}}$, un point maximal
$x'$ de $(T_\alpha)_{y'}$ est associé à $G_{y'}$ (« EGA IV 12.2 »), donc dans
$U_{y'}$. Le contre-exemple montre ce qui arrive quand (a) manque : $X \to Y$
birationnel dont une fibre se relève en une composante verticale, $G =
\mathcal{O}_X$ ; la fibre de $T = X$ y est de dimension $1$ au lieu de
$0$\footnote{Avec un croquis : une courbe $X$ au-dessus d'une droite $Y$.
La première ligne du contre-exemple, « (l'ouvert i-fini est $\simeq$) », est
une lecture peu sûre.}.

\subsection*{Les remarques et le « Vrai théorème » (page 15)}

\begin{itemize}
  \item \emph{1.4.2, remarque 1.} Les conditions (b) et (c) de 1.4 équivalent
  à : le problème a une solution après tout changement de base $Y' \to Y$ avec
  $Y'$ artinien.
  \item \emph{1.4.3, remarque 2.} La condition (a) résulte de l'existence
  d'une solution après tout changement de base $Y' \to Y$ avec $Y'$ un trait.
\end{itemize}

D'où ce que la page appelle le \emph{Vrai théorème}, et, en marge, l'«
Énoncé final » : sous les hypothèses préliminaires de 1.4, le problème a une
solution si et seulement s'il en a une après tout changement de base
$Y' \to Y$ avec $Y'$ artinien ou un trait.
Le dossier ne le démontre pas\footnote{Une note encerclée, reliée à « trait »,
porte « inutile si les $\mathcal{O}_{Y,y}$ réduits ». Il n'est pas clair
qu'elle porte sur les traits plutôt que sur les bases artiniennes : le
corollaire 1.6, qui suppose justement les $\mathcal{O}_{Y,y}$ réduits, ne garde
que les traits. Nous ne tranchons pas.}.

La seconde moitié de la page, barrée, commence la démonstration du
corollaire 1.5 : par le vrai théorème, on se ramène au cas où la base est un
trait ; si $x \in \mathrm{Ass}\, G_U$ est au-dessus du point générique $y'$ et
$T = \overline{\{x\}}$ rencontre la fibre spéciale, les composantes de
$T_y$ sont de dimension $\geqslant \dim T_{y'} = d$, donc sont des
composantes de $X_y$.

\pagerange{16}{16}
\subsection*{La mise au net d'une autre main (page 16)}

La page 16 n'est pas de la main de Grothendieck. Elle met au net l'énoncé
1.4, sous le nom de « Lemme 1.4' », avec les notations de 1.1, $X$ et $Y$
localement noethériens, $G$ cohérent, $Z$ l'espace sous-jacent à $X - U$ et
$y \in f(Z)$. La condition sur les fibres formelles y devient une condition
(d) aux points de $Z_y$ ; la condition (a) y est la forme par densité donnée
plus haut, demandée aussi pour tout $X'$ étale sur $X$ ; (b) et (c) sont
celles que nous avons écrites. Sa conclusion est locale en $y$ : $G$ est plat
sur $Y$ en tout point de $X_y$, et $U_y$ contient $\mathrm{Ass}\, G_y$ ; si les
conditions valent en tout point de $f(Z)$, le problème a une solution
cohérente\footnote{Nous ne donnons de cette page que ce que la note de la
transcription en dit. Aucune intervention de Grothendieck n'y est sûre.}.

\pagerange{18}{19}
\subsection*{L'ouvert universellement dense, et le rôle de (b) (pages 18 et 19)}

La page 18 reprend la condition de 1.1 pour un sous-schéma fermé $Z$ de
$X$, plat sur $Y$, et $Z_U = Z \cap U$ : $Z_U$ doit être universellement
schématiquement dense dans $Z$ relativement à $Y$ ; si $Z$ est localement de
présentation finie sur $Y$, ou si $Y$ et $Z$ sont noethériens, cela équivaut à
$\mathrm{Ass}\, Z_y \subset U$ pour tout $y$, avec un renvoi à « EGA IV
11.12\,? »\footnote{Le mot lu « univoque » devant « condition » est peu sûr. Le
point d'interrogation après le renvoi est de la page. L'équivalence est celle
que nous avons citée pour 1.1.}.

La page 19 donne la clef de l'étape 4 de la démonstration de 1.4. Posons
$A = \widehat{\mathcal{O}}_{Y,y}$. Pour tout $n$, il existe une partie
\emph{finie} $J$ de l'ensemble d'indices telle que
\[
  \mathfrak{A}_n := \mathrm{Ker}\Bigl(A \to \prod_{\alpha \in J}
  \mathcal{O}_{Y_\alpha, y_\alpha}\Bigr) \subset \widehat{\mathfrak{m}}_y^{\,n} ,
\]
et l'on peut supposer les $Y_i$ artiniens. La page n'en dit pas plus ; voici
comment ces deux faits s'enchaînent\footnote{L'enchaînement est de nous. Il
n'utilise que des énoncés du dossier : la page 19, le lemme 1.3, la condition
(c).}.

\begin{itemize}
  \item L'existence de $J$ est le lemme de Chevalley : $A$ est local
  noethérien complet, donc linéairement compact, et une famille filtrante
  décroissante d'idéaux fermés d'intersection nulle finit par entrer dans
  chaque $\widehat{\mathfrak{m}}_y^{\,n}$.
  \item On peut supposer les $Y_i$ artiniens en les remplaçant par leurs
  tronqués $\mathrm{Spec}\bigl(\mathcal{O}_{Y_i}/\mathfrak{m}_i^{\,k}\bigr)$ :
  l'intersection des noyaux reste nulle (Krull dans chaque
  $\mathcal{O}_{Y_i}$), et la finitude de (c) se conserve.
  \item Alors $A/\mathfrak{A}_n$ est artinien, et
  $\coprod_{\alpha \in J} Y_\alpha \to \mathrm{Spec}(A/\mathfrak{A}_n)$ est
  surjectif et injectif sur les fonctions : c'est un épimorphisme, après
  lequel le problème a une solution. Le lemme 1.3, dont (c) fournit les
  hypothèses, donne une solution sur $\mathrm{Spec}(A/\mathfrak{A}_n)$, donc
  sur $\mathrm{Spec}(A/\widehat{\mathfrak{m}}_y^{\,n})$ par changement de base.
\end{itemize}

C'est la remarque 1.4.2 dans un sens : (b) et (c) donnent des solutions sur
tous les voisinages infinitésimaux du point $y$, et l'étape 4 passe à la
limite sur ceux-là. La page porte aussi des croquis de la réduction (les
schémas $X \leftarrow X' \leftarrow X''$ au-dessus de $Y \leftarrow Y'
\leftarrow Y''$, le premier morphisme fpqc) et six croquis de surfaces et de
courbes qui ne sont pas reproduits\footnote{Sur la page, une double flèche
entre $Y'$ et $Y''$ et une triple entre $Y''$ et $Y'''$ : c'est le début du
complexe de descente.}.

\pagerange{5}{6}
\section*{IV. Les corollaires 1.5 à 1.7 (pages 5 et 6)}

\emph{Corollaire 1.5.} Supposons $f$ localement de type fini, $F$ et $G$
cohérents, et (a') les fibres $(G_U)_y$ vérifiant $(S_1)$ et équidimensionnelles
d'une même dimension $d$ ; supposons $U_y$ dense dans $X_y$ pour tout $y$, et
les anneaux locaux $\mathcal{O}_{X,x}$, $x \in X - U$, à fibres formelles
géométriquement normales. Alors la condition (b) de 1.4 est nécessaire et
suffisante pour l'existence d'une solution.

Toute cette moitié de page est barrée d'une grande croix au
crayon\footnote{Nous rapportons donc l'énoncé sans l'asserter. La page ajoute
qu'il suffit que $\mathcal{O}_{Y,f(x)}$ ait des fibres formelles
géométriquement normales : c'est une instance de la permanence des propriétés
des fibres formelles par passage aux algèbres de type fini (EGA IV, § 7.4,
cité de mémoire, sans vérification des hypothèses exactes). Une note de marge,
encerclée et barrée, lit « S-M $G = X$ ? ».}. Son idée est claire : sous (a')
la condition (a) devient automatique (c'est la fin de la page 15), et il ne
reste que (b).

\emph{Critère valuatif.}

\emph{Corollaire 1.6.} Plaçons-nous sous les hypothèses préliminaires de 1.4,
y compris celle sur les fibres formelles. Supposons de plus que, pour tout
$x \in X - U$ et $y = f(x)$, l'anneau $\mathcal{O}_{Y,y}$ soit réduit, et que
pour tout trait $Y'$ fini sur $\mathrm{Spec}\, \mathcal{O}_{Y,y}$ le problème
ait une solution après le changement de base $Y' \to Y$. Alors le problème a
une solution\footnote{Deux notes au crayon dans la marge : « 1.5\,? \ldots,
mais sans hypothèse du type de 1.4 a) », et « Fait double emploi avec 1.7,
sauf qu'on a fait d'hyp.\ « de type fini » ».}.

\emph{Corollaire 1.7.} Soient $Y$ localement noethérien, $F$ et $G_U$
cohérents. On suppose :
\begin{itemize}
  \item[a)] pour tout $x \in X - U$ et $y = f(x)$, $\mathcal{O}_{Y,y}$ est
  réduit, à fibres formelles géométriquement normales (c'est le cas, par
  exemple, s'il est essentiellement de type fini sur $\mathbb{Z}$)\footnote{La
  parenthèse est en partie illisible (« cf.\ \ldots\ de type fini sur
  $\mathrm{Spec}\, \mathbb{Z}$ »). L'exemple est exact : un anneau
  essentiellement de type fini sur $\mathbb{Z}$ est excellent, ses fibres
  formelles sont géométriquement régulières, a fortiori normales.} ;
  \item[b)] pour un tel $y$ et tout trait $Y'$ fini sur
  $\mathrm{Spec}\, \mathcal{O}_{Y,y}$, le problème a une solution après
  $Y' \to Y$.
\end{itemize}
Alors le problème posé a une solution (page 6). La page écrit « Explicitez »
après b), et une troisième condition c) est barrée.

Une note de marge dit pourquoi le critère est maniable : dans la condition b),
la platitude est automatique, et seule compte la question des points
associés\footnote{« NB Dans cette condition b) l'hypothèse de platitude devient
triviale, seule la question des $\mathrm{Ass}\, G'_{y'}$ \ldots\ est
sérieuse ». C'est exact (de nous) : sur un trait, plat veut dire sans torsion ;
or $G \subset i_*(G_U)$, et $i_*(G_U)$ est sans torsion puisque $G_U$
l'est.}. C'est l'analogue, pour le problème 1.1, du critère valuatif de
platitude d'EGA IV, § 11.8, cité de mémoire.

La page 6 se termine sur : « Nous allons maintenant utiliser le \emph{lemme
de Hironaka}. »

\pagerange{7}{7}
\section*{V. Le théorème 1.8 et le lemme de Hironaka (pages 7, 9, 10 et 12)}

\subsection*{L'énoncé de la page 7}

\emph{Théorème 1.8.} Soient $f \colon X \to Y$ localement de présentation
finie, $x \in X$ et $y = f(x)$. On suppose :
\begin{itemize}
  \item[a)] $Y$ réduit en $y$ ;
  \item[b)] $f$ universellement ouvert aux générisations maximales de $x$ dans
  $X_y$ ;
  \item[c)] pour toute générisation $y'$ de $y$ et toute générisation $x'$ de
  $x$ dans $X_{y'}$, $\dim_{x'} X_{y'} = \dim_x X_y$ ;
  \item[d)] $X_y$ réduit aux générisations maximales de $x$ dans $X_y$ ;
  \item[e)] $(X_y)_{\mathrm{red}}$ géométriquement normal au voisinage de $x$.
\end{itemize}
Les conditions d) et e) équivalent à l'existence d'un voisinage ouvert $V$ de
$x$ dans $X_y$ tel que $V \otimes_{k(y)} \overline{k(y)}$ soit réduit aux
points génériques et que sa partie réduite soit normale.

\emph{Alors} il existe un voisinage ouvert $U$ de $x$ tel que
$U_{\mathrm{red}} \to Y$ soit normal, c'est-à-dire plat à fibres
géométriquement normales\footnote{La conclusion de la page 7 est barrée et lue
en partie seulement ; « normal » y est d'abord biffé. Nous donnons la
conclusion de la remise au net de la page 10, qui dit « normal », i.e.\
« \emph{plat} à fibres géom.\ normales ». Une note de marge propose de
remplacer b) par « $f$ universellement ouvert au voisinage de $x$ », une autre
remarque que, si $X$ et $Y$ sont localement noethériens, on doit pouvoir se
passer de la présentation finie moyennant des hypothèses d'excellence.}.

Un « Cor (2) », entièrement encadré de traits obliques, cherche des conditions
équivalentes en posant $Y_0 = Y_{\mathrm{red}}$, $X_0 = X \times_Y Y_0$ :
(i bis) $f$ universellement ouvert et équidimensionnel, à fibres
géométriquement réduites et géométriquement normales ; (ii) un ouvert $V$ de
$X_0$, dense dans chaque fibre, sur lequel $V \to Y_0$ est séparable, les
$(X_{0,y})_{\mathrm{red}}$ géométriquement normaux ; (iiii bis) un
sous-schéma $Z_0$ de $X_0$, localement de présentation finie sur $Y_0$, normal,
coïncidant avec $X_0$ sur un ouvert schématiquement dense\footnote{La
numérotation (i), (i bis), (ii), (iii), (iiii bis) est celle de la page ;
(iii) et une partie de (iiii bis) sont illisibles. Ce corollaire est rapporté,
non asserté.}.

\pagerange{9}{9}
\subsection*{La démonstration esquissée (page 9)}

La page 8 est le feuillet d'EGA IV dont il est question plus bas ; la
démonstration est à la page 9\footnote{La page 9 porte « 7 » dans sa
numérotation, surchargeant un « 8 » : le feuillet dactylographié a été
glissé là.}.

\emph{Reformulation.} La conclusion (ii) du Cor (2) dit aussi : il existe un
sous-schéma $U_0$ de $U$, de présentation finie sur $Y_0$, tel que $U_0 \to U$
soit surjectif et $U_0 \to Y$ normal --- ce qui force $U_0$ réduit, donc
$U_0 = U_{\mathrm{red}}$ --- et un ouvert $V$ de $U_0$, schématiquement dense
relativement à $Y$, tel que $V \to X$ soit une immersion ouverte. Sous cette
forme, les techniques de passage à la limite d'EGA IV ramènent au cas où $Y$
est le spectre d'un anneau local essentiellement de type fini sur
$\mathbb{Z}$\footnote{La lettre notée $U$ est une capitale surchargée, entre
$U$ et $V$.}.

\emph{Réduction à la fibre.} Par constructibilité, passage à la limite et le
lemme de Hironaka, on se ramène à des hypothèses sur la fibre (EGA IV 16.1
est cité pour la dimension). La condition c) donne que $f$ est universellement
ouvert et équidimensionnel au voisinage de $x$. On peut supposer $X$ réduit,
et il s'agit de prouver que $X \to Y$ est normal.

\emph{L'ouvert séparable.} Soit $V$ l'ensemble des points où $f$ est plat à
fibres géométriquement réduites (« séparable »). Il est ouvert (EGA IV, § 12),
il contient les points de $X_y$ où $X_y$ est géométriquement réduit (EGA IV
15.2.3, qui sert ici parce que $Y$ est réduit et $f$ universellement ouvert),
et il est donc dense dans chaque fibre par d) et e). Le reste, au bas de la
page, est encadré et barré, sauf deux mots : « (par morceaux\,!) \ldots\ c'est
le cas de Hironaka, OK »\footnote{Deux lignes y sont écrites tête-bêche et ne
sont pas lues. En marge : « Lemme à dégager », souligné deux fois ; « On peut
procéder par récurrence sur $\dim Y$ », barré ; et, au crayon, « Le critère
valuatif nous ramène au cas où $Y$ est le spectre d'un anneau de val.\
discrète ».}.

\pagerange{10}{10}
\subsection*{La remise au net (page 10)}

La page 10 récrit l'énoncé pour $f$ de type fini et $Y$ noethérien :
$\mathcal{O}_{Y,y}$ et $\mathcal{O}_{X,x}$ réduits, $f$ universellement ouvert
(« EGA IV 14, 15 »), les composantes irréductibles des fibres de même
dimension $d$ au voisinage de $x$, $X_y$ réduit aux générisations maximales de
$x$, $(X_y)_{\mathrm{red}}$ géométriquement normal sur $k(y)$ ; alors $x$ a un
voisinage ouvert $U$ tel que $U \to Y$ soit plat à fibres géométriquement
normales\footnote{La page écrit « voisinage ouvert $U$ de $X$ » pour « de
$x$ ».}.

La démonstration est donnée lorsque $\mathcal{O}_{Y,y}$ est excellent --- il
suffit, précise la page, qu'il soit à fibres formelles géométriquement
normales, ce qui est l'hypothèse du § 1. a) Les hypothèses restent vraies au
voisinage de $x$ (renvoi « EGA IV 9 », avec un croquis de changement de base
$Y' \to Y$). Puis : l'ensemble $V$ des points où $X/Y$ est plat à fibres
séparables est ouvert (EGA IV, § 12), contient les points de $X_y$ où $X_y$ est
réduit (EGA IV 15.2.3), donc est dense dans $X_y$, et l'on conclut par
équidimensionnalité qu'il est dense dans chaque fibre près de $x$.

\pagerange{12}{12}
\subsection*{La fin de la preuve (page 12)}

b) On suppose les hypothèses satisfaites en tout point de $X$. Soit $U$
l'ouvert où $f$ est séparable : il est dense dans chaque fibre, et il reste à
prouver que $X = U$ et que $X$ est plat sur $Y$, autrement dit que $X$ est
l'adhérence schématique de $U$ et qu'elle est plate.

\emph{Normalité des fibres.} On se ramène à prouver $X_y$ séparable sur
$k(y)$. Soit $y'$ une générisation maximale de $y$, et un trait $Y' \to Y$
reliant $y'$ à $y$. Après ce changement de base, la fibre spéciale de
l'adhérence schématique de la fibre générique est réduite ; le lemme de
Hironaka s'applique, elle est normale, donc $X_y$ aussi\footnote{Les indices
$y'_0$ de la page sont peu sûrs, et la phrase « donc $\overline{(\ )}$ (adh.\
sch.)\ $= X'_{y'_0}$ » est écrite serrée. Le lemme de Hironaka n'est pas énoncé
dans le dossier ; sous sa forme usuelle (EGA IV 5.12.8, cité de mémoire), il
dit qu'un schéma plat sur un trait, à fibre générique normale, dont la fibre
spéciale est réduite à ses points génériques et de partie réduite normale, a
une fibre spéciale normale.}.

\emph{Platitude.} c) Pour prouver la platitude de $X/Y$, on applique le
« critère valuatif de platitude d'une adhérence schématique », c'est-à-dire
les corollaires 1.6 et 1.7 avec $F = \mathcal{O}_X$ et $G_U = \mathcal{O}_U$ :
c'est là que la première partie du dossier sert. Le reste du feuillet est
blanc.

\pagerange{8}{8}
\section*{VI. Un feuillet d'EGA IV (page 8)}

La page 8 est une page dactylographiée du tapuscrit d'EGA IV, § 21 : la fin
d'une démonstration, puis le n° 21.10, « Factorialité des anneaux locaux
réguliers », et le début de la preuve du théorème 21.10.1 (un anneau local
régulier est factoriel). Le texte est publié et n'est pas restitué ici. Sa
main y ajoute, en marge, « Créditer Kaplansky de la démonstration ici », et
au-dessus du théorème « (Auslander-Buchsbaum) » ; il corrige un renvoi
(« (5.11.6) » devient « I 9.4.5 »), biffe une proposition et numérote une
résolution (21.10.1.1). Le feuillet n'est pas un brouillon de ce dossier ; sa
présence est compatible avec la datation du groupe, mais ne date
rien\footnote{Le numéro de page dactylographié, IV-1131, est complété à la
main et lu avec doute.}.

\pagerange{11}{11}
\section*{VII. Le site des schémas finis et les faisceaux formels (pages 11, 13, 17 et 20)}

\subsection*{Quasi-cohérence sur les schémas finis (page 11)}

Soient $k$ un corps algébriquement clos, $S = X = \mathrm{Spec}\, k$. La page
fixe trois classes : $S_0$, les schémas étales de type fini sur $X$ ;
$\mathfrak{S}$, les schémas de type fini sur $S$ ; $\mathfrak{M}$, les
immersions ouvertes\footnote{La glose de $S_0$, lue « revêtements locaux »,
est peu sûre ; les lettres $\mathfrak{S}$, $\mathfrak{M}$ sont des capitales de
forme gothique.}.

La catégorie $\mathcal{C}$ est celle des schémas finis sur $k$, avec la
topologie de Zariski. Un schéma fini est réunion disjointe de spectres
d'anneaux locaux artiniens, et ses recouvrements de Zariski sont ses
décompositions en parties ouvertes et fermées ; un faisceau sur $\mathcal{C}$
est donc un foncteur $\Lambda \mapsto F(\Lambda)$ sur les $k$-algèbres finies
qui transforme produits finis en produits\footnote{La justification est de
nous. La page ajoute, peu lisiblement, « \& connexes » : si l'on se borne aux
schémas finis connexes, la condition de faisceau disparaît.}. Les faisceaux
quasi-cohérents sur $\mathrm{Spec}\, \Lambda$ sont les
$\Lambda' \mapsto M \otimes_\Lambda \Lambda'$, pour $M$ un $\Lambda$-module ;
ceux de présentation finie correspondent aux $M$ de type fini.

Mais le noyau d'un morphisme de faisceaux quasi-cohérents n'est pas
quasi-cohérent. Pour $u \colon \mathcal{O}_\Lambda^{\,n} \to
\mathcal{O}_\Lambda$ défini par $u_1, \ldots, u_n \in \Lambda$, le faisceau
$\Lambda' \mapsto \mathrm{Ker}\bigl(u_{\Lambda'} \colon \Lambda'^{\,n} \to
\Lambda'\bigr)$ n'est pas en général $\Lambda' \mapsto (\mathrm{Ker}\, u)
\otimes_\Lambda \Lambda'$ : « manque de platitude ». Exemple, de nous :
$\Lambda = k[\varepsilon]$, $n = 1$, $u = \varepsilon$ ; $\mathrm{Ker}\, u =
\varepsilon \Lambda$, mais pour $\Lambda' = k$ le noyau de $u_k = 0$ est $k$
tout entier, alors que $(\varepsilon\Lambda) \otimes_\Lambda k \to k$ est nul.

Le bas de la page porte des croquis au crayon qui reprennent les pages 3 et 4
(« $U \cap T_y$ dense dans $T_y$ »), sans rapport avec ce qui précède.

\pagerange{13}{13}
\subsection*{Faisceaux formels (pages 13 et 17)}

Soit $\mathcal{O}$ le faisceau d'anneaux standard sur $\mathcal{C}$ et $F$ un
$\mathcal{O}$-module quasi-cohérent. Alors $F$ induit sur chaque
$\mathcal{M}_{/T}$ un faisceau quasi-cohérent, de façon cartésienne en l'objet
variable $T_0 \to T$\footnote{Note de marge : « à vérifier ». Les petits $X$
surmontant les flèches sont peu sûrs, et une grande partie du second
paragraphe de la page 13 est illisible.}. On obtient ainsi un « faisceau
formel » sur le complété formel $\widehat{Z}$ d'un schéma $Z$ le long de
$X$, pour $\delta \colon X \to Z$, et l'on annonce
\[
  H^*(\mathcal{C}_{/\delta}, F) \simeq H^*(\widehat{Z}, F|_{\widehat{Z}}) ,
\]
le second membre étant calculé dans la catégorie des systèmes projectifs
adiques sur $\widehat{Z}$. La page marque la formule « à vérifier » ; le dossier
ne la vérifie pas.

\pagerange{17}{17}
\subsection*{Les groupes explicités (page 17)}

La page 17 explicite ces groupes. Si $\delta$ est dans $\mathcal{N}$,
c'est-à-dire $Z \in \mathrm{Ob}\, \mathfrak{S}$, la catégorie
$\mathcal{C}_{/\delta}$ des carrés cartésiens au-dessus de $X \to Z$ est
équivalente à $\mathcal{M}_{/Z}$ munie de la topologie induite, et l'on trouve,
pour les voisinages infinitésimaux $X^{(i)}_n$ supposés dans $\mathfrak{S}$, la
suite exacte
\[
  0 \to {\varprojlim_n}^{1} H^{q-1}(X^{(i)}_n, F) \to H^q(X^{(i)}, F) \to
  \varprojlim_n H^q(X^{(i)}_n, F) \to 0
\]
(les groupes de droite étant pris pour la topologie de $\mathcal{M}$ sur
$\mathcal{C}$)\footnote{L'exposant du premier terme se lit $q$ sur la page,
dans une formule encadrée et surchargée ; c'est $q - 1$ qui est juste, et c'est
ce qu'écrit la page 20.}.

\pagerange{20}{20}
\subsection*{La suite de Milnor (page 20)}

La page 20 démontre cette suite dans un cadre général. Soient $S_0 \subset
S_1 \subset \cdots$ et $S_\infty$ leur limite, $F$ un faisceau, et
$C^{\cdot}(F)$ une résolution injective de $F$ ; ses restrictions sont des
résolutions injectives des $F_{S_i}$, et
\[
  \Gamma\bigl(S_\infty, C^{\cdot}(F)\bigr) = \varprojlim_i
  \Gamma\bigl(S_i, C^{\cdot}(F)\bigr) ,
\]
le système projectif étant strict, c'est-à-dire à flèches de transition
surjectives. Il en résulte que $H^n(S_\infty, F) \to \varprojlim H^n(S_i, F)$
est surjectif, et bijectif si le système des $H^{n-1}(S_i, F)$ vérifie la
condition de Mittag-Leffler. En général, on a une suite spectrale
$E_2^{pq} = {\varprojlim_i}^{(p)} H^q(S_i, F) \Rightarrow H^{p+q}(S_\infty,
F)$, et comme ${\varprojlim}^{(p)} = 0$ pour $p \geqslant 2$, des suites exactes
\[
  0 \to {\varprojlim_i}^{1} H^{n-1}(S_i, F) \to H^n(S_\infty, F) \to
  \varprojlim_i H^n(S_i, F) \to 0 .
\]
C'est la suite exacte de Milnor\footnote{L'annulation de
${\varprojlim}^{(p)}$ pour $p \geqslant 2$ suppose le système indexé par
$\mathbb{N}$, ou par un ensemble dénombrable ; la page ne le dit pas, mais ses
indices sont des entiers. La page écrit $C^{\cdot}(F)_{S_\infty}$ et
$C^{\cdot}(F)_{S_i}$ pour les restrictions ; le nom de Milnor n'y figure
pas.}. La page s'arrête sur « Dans le cas qui nous occupe, il vient : ».

\pagerange{21}{22}
\section*{VIII. Foncteurs sur les schémas finis (pages 21 à 23)}

La chemise de la page 21 porte « Foncteurs sur morphismes non ramifiés, étales
\ldots »\footnote{Encre brune, en haut à droite d'un feuillet par ailleurs
blanc. Un mot est biffé après « Morphismes » ; la note de marge de la page
22, « Avec morphismes non ramifiés », appuie les lectures « non » et
« ramifiés ».}.

\subsection*{Le lemme (page 22)}

\emph{Lemme.} Soient $S$ un préschéma localement noethérien, $\mathcal{C}$ la
catégorie des préschémas finis sur $S$, et $F$ un foncteur contravariant de
$\mathcal{C}$ dans les ensembles. On suppose :
\begin{itemize}
  \item[(i)] pour tout $X \in \mathrm{Ob}\, \mathcal{C}$, l'application
  $F(X) \to F(X_{\mathrm{red}})$ est injective ;
  \item[(ii)] pour tout épimorphisme \emph{effectif} $X \to Y$ de
  $\mathcal{C}$, la suite $F(Y) \to F(X) \rightrightarrows F(X \times_Y X)$
  est exacte.
\end{itemize}
Alors la suite $F(Y) \to F(X) \rightrightarrows F(X \times_Y X)$ est exacte
pour \emph{tout} épimorphisme $X \to Y$ de $\mathcal{C}$\footnote{La condition
(i) primitive est biffée ; la condition actuelle est ajoutée en interligne,
« injectif » à l'encre brune. « Exacte » s'entend au sens des ensembles :
$F(Y)$ s'identifie à l'égalisateur des deux flèches.}.

\emph{Démonstration.} Elle repose sur le fait, que la page utilise sans le
démontrer, que tout épimorphisme de $\mathcal{C}$ est composé d'un nombre fini
d'épimorphismes effectifs\footnote{Le dossier 128 du même groupe étudie, pages
12 à 18, des extensions finies $A \subset A'$ d'anneaux locaux et l'exactitude
de $A \to A' \rightrightarrows A' \otimes_A A'$ ; le rapprochement est de nous,
et nous ne savons pas si ce fait est vrai en toute généralité.}. Par (ii),
chaque facteur induit une injection, donc $F(Y) \to F(X)$ est injectif pour
tout épimorphisme. On raisonne par récurrence sur le nombre $n$ de facteurs ;
pour $n = 1$, c'est (ii). Si $n > 1$, on écrit
\[
  X \xrightarrow{\ u'\ } X' \xrightarrow{\ u''\ } Y ,
\]
$u'$ composé de $n - 1$ épimorphismes effectifs, $u''$ effectif. Soit
$\xi \in F(X)$ dont les deux images dans $F(X \times_Y X)$ sont égales.
\emph{A fortiori} ses deux images dans $F(X \times_{X'} X)$ le sont, et par
récurrence $\xi$ provient d'un $\xi' \in F(X')$. Il reste à voir que les deux
images de $\xi'$ dans $F(X' \times_Y X')$ coïncident ; comme elles s'envoient
sur les deux images de $\xi$, il suffit que
\[
  F(X' \times_Y X') \to F(X \times_Y X)
\]
soit injectif. La page s'arrête là\footnote{La dernière ligne, serrée au bas du
feuillet, s'arrête sur « lui-même conséquence \ldots\ ceci ». La page écrit
d'abord $X'_1$, puis $X'$. Le diagramme de la marge, que nous ne reproduisons
pas, montre $F(Y) \to F(X') \to F(X)$ au-dessus de
$F(X' \times_Y X') \to F(X \times_Y X) \to F(X \times_{X'} X)$.}. C'est le
corollaire de la page 23 qui fournit ce qui manque : $X \times_Y X \to X'
\times_Y X'$ est surjectif, comme produit de deux morphismes surjectifs, et
son argument n'utilise que l'injectivité déjà acquise, si bien qu'il n'y a pas
de cercle\footnote{Cet enchaînement est de nous.}.

\pagerange{23}{23}
\subsection*{Le corollaire et les sous-schémas abéliens (page 23)}

\emph{Corollaire.} Sous (i) et (ii), si $u \colon X \to Y$ est
\emph{surjectif}, alors $F(Y) \to F(X)$ est \emph{injectif}. En effet, par (i)
et le carré
\[
  \begin{array}{ccc}
    F(Y) & \longrightarrow & F(X) \\
    \downarrow & & \downarrow \\
    F(Y_{\mathrm{red}}) & \longrightarrow & F(X_{\mathrm{red}})
  \end{array}
\]
on se ramène à $X$ et $Y$ réduits ; alors $u$ est un épimorphisme de
$\mathcal{C}$, et l'on a déjà vu que $F(Y) \to F(X)$ est
injectif\footnote{Pourquoi $u$ est alors un épimorphisme (de nous) : $Y$ étant
réduit et $u$ surjectif, $\mathcal{O}_Y \to u_*\mathcal{O}_X$ est injectif ;
comme les objets de $\mathcal{C}$ sont affines sur $S$, deux morphismes de $Y$
vers un objet de $\mathcal{C}$ qui coïncident sur $X$ coïncident. Le carré est
en marge, à l'encre brune.}.

\emph{Proposition (cas particulier).} Soient $S$ localement noethérien, $A$ un
schéma abélien sur $S$, et pour $S'$ fini sur $S$, $F(S')$ l'ensemble des
sous-schémas abéliens de $A' = A \times_S S'$. Si $S' \to S$ est un
épimorphisme, la suite
\[
  F(S) \to F(S') \rightrightarrows F(S'') , \qquad S'' = S' \times_S S' ,
\]
est exacte\footnote{« abélien » et « abéliens » sont des lectures peu sûres.
La page écrit « pour tout $S'$ sur $S$ » ; le lemme ne s'applique qu'à
$S'$ fini sur $S$, et nous l'ajoutons. Elle écrit aussi
« $0 \to F(S) \to \cdots$ », que nous lisons comme l'injectivité de
$F(S) \to F(S')$.}. La condition (i) du lemme résulte de la rigidité : un
sous-schéma abélien est déterminé par sa réduction ; la condition (ii) résulte
de la théorie de la descente.

La page commence un « Corollaire » (d'abord « Proposition »), dont la
parenthèse porte un nom biffé et un autre, illisibles, et annonce un « Énoncé
analogue pour les \ldots\ abéliens ». L'énoncé s'arrête là, et le dossier
avec lui.

\end{document}
